\documentclass[11pt,a4paper]{moderncv}
\moderncvtheme[blue]{classic}

\usepackage[utf8]{inputenc}
\usepackage[top=0.5cm, bottom=0.5cm, left=0.5cm, right=0.5cm]{geometry}
\usepackage{graphicx}
\usepackage{multibib}
\usepackage{fontawesome5} % para íconos de redes sociales

% Datos personales
\firstname{David}
\familyname{Valenciano}
\title{Senior Front-End Engineer}
\mobile{+34 600 80 90 24}
\email{david.valenciano.esteban@gmail.com}
\homepage{davalest.github.io/resume}
\extrainfo{linkedin.com/in/david-valenciano \quad github.com/davalest}
%\photo[50pt][0.4pt]{photo.jpg}
% Handles corregidos: los que estaban comentados apuntaban a perfiles inexistentes.
%\social[linkedin]{david-valenciano}
%\social[github]{davalest}

% Otros ajustes
\newcites{book,misc}{{Books},{Others}}
\nopagenumbers{}

\begin{document}
    \maketitle

    \section{Perfil}
    \cvitem{}{Front-end engineer con más de 10 años de experiencia, 9 de ellos con React.
    Construyo sistemas de componentes que escalan, código tipado que caza los problemas antes
    de que lleguen a QA, y aplicaciones web y multiplataforma que aguantan en producción.
    Ahora en Docline (healthtech), en la plataforma de consulta médica online y en la librería
    de componentes que comparten sus productos.}

    \section{Experiencia Profesional}
    \cventry{Ene 2024 -- Actualidad}{Senior Front-End Engineer}{Docline (Aplicaciones de Salud, S.L.)}{}{}{
        \begin{itemize}
            \item Uno de los dos front-end que cubren los cinco productos de la empresa: \textit{VideoConsultation}, \textit{Portal}, \textit{Backoffice}, \textit{Docline Care} y \textit{Docline Components}.
            \item Responsable de \textit{Docline Components}, la librería sobre la que se construyen todos esos productos. React, Styled-Components y TypeScript.
            \item Defino la arquitectura de front y los estándares de código, y llevo el code review y la mentoría del equipo.
            \item Desarrollé \textit{Docline Care} en solitario con React Native y Expo, y me encargué de su publicación en App Store y Google Play.
            \item Lidero la adopción de IA en Docline desde marzo de 2026: introduje una metodología de desarrollo agéntico con IA (BMAD), la usé para construir el back-office de la empresa de principio a fin, y ahora imparto mentorías y workshops para extenderla al resto del equipo.
        \end{itemize}
    }

    \cventry{Ago 2022 -- Nov 2023}{Front-End Engineer}{Shopery, S.L.}{}{}{
        \begin{itemize}
            \item \textit{Mustang} y \textit{YouPop}: tiendas con Next.js, React.js, TypeScript y SCSS, en producción en España, México y República Dominicana. Equipo de cinco front-end.
        \end{itemize}
    }

    \cventry{Sep 2021 -- Ago 2022}{Front-End Engineer}{Stack and Vault, S.L.}{}{}{
        \begin{itemize}
            \item \textit{Wotoch}: app web con Svelte y TypeScript.
            \item \textit{Voting Manager}: app híbrida con React.js, React Native, HTML5, SCSS y TypeScript.
            \item \textit{VR Manager} e \textit{ImmoScout24}: apps web con React.js, HTML5, SCSS, TypeScript y Redux.
        \end{itemize}
    }

    \cventry{Oct 2019 -- Sep 2021}{Analista Programador}{Capitole Consulting}{}{}{
        \begin{itemize}
            \item \textit{Verisure OWA}: app web con React.js, Redux, GraphQL, HTML5 y SCSS.
        \end{itemize}
    }

    \cventry{Feb 2017 -- Oct 2019}{Front-End Developer}{QUO Health, S.L.}{}{}{
        \begin{itemize}
            \item \textit{gluQUO Diabetes Academy}: app híbrida y web con React.js, React Native, HTML5, SCSS, Firebase.
            \item \textit{O.K Transport}: app web con React.js, HTML5, CSS3, Bootstrap, Node.js, MongoDB.
            \item \textit{gluQUO PRO}, \textit{Optima cc web}, \textit{gluQUO web}, \textit{QUO Health web}: apps web con React.js, HTML5, CSS3, Bootstrap.
        \end{itemize}
    }

    \cventry{2016}{Front-End Developer, freelance}{}{}{}{
        \begin{itemize}
            \item \textit{MarVal S.L.} y \textit{La Parada S.L.}: diseño, maquetación y desarrollo web con HTML5, CSS3, Bootstrap y JQuery.
        \end{itemize}
    }

    \cventry{Sep 2015 -- Jul 2016}{Soporte IT (CAU)}{HM Hospitales (Montepríncipe)}{}{}{
        \begin{itemize}
            \item Mantenimiento, gestión de incidencias e implementación de infraestructura.
        \end{itemize}
    }

    \section{Stack}
    \cvitem{}{
        \begin{tabular}{p{0.26\linewidth} p{0.68\linewidth}}
            \textbf{A diario} & React, TypeScript, React Native, Expo, Styled-Components, SCSS \\
            \textbf{En producción} & Next.js, Svelte, Redux, GraphQL \\
            \textbf{Backend (2017--2019)} & Node.js, Express, Firebase, MongoDB \\
        \end{tabular}
    }

    \section{Formación Académica}
    \cventry{2024}{Técnico Superior en Desarrollo de Aplicaciones Multiplataforma}{Ilerna}{}{Online}{}
    \cventry{2020}{Técnico Superior en Desarrollo de Aplicaciones Web}{Ilerna}{}{Online}{}
    \cventry{2017}{Técnico Superior en Administración de Sistemas Informáticos y Redes}{IES Lázaro Cárdenas}{}{Collado Villalba}{}
    \cventry{2015}{Técnico en Sistemas Microinformáticos y Redes}{IES Infanta Elena}{}{Galapagar}{}

    \section{Información Adicional}
    \cvline{Idiomas}{Castellano (nativo), Inglés (fluido, tras vivir en Dublín)}

\end{document}
